% TOP_PROBLEM: {"id": 3047, "title": "Dividing up the unrestricted partitions", "category_id": 2, "category": {"id": 2, "name": "combinatorics", "display_name": "Combinatorics", "description": "Counting problems, graph theory, discrete structures.", "slug": "combinatorics", "order_index": 2, "created_at": "2026-07-31T15:26:25.671Z"}, "status": "open", "problem_number": "OPG-37222", "set_id": 12}
% FIRST_POSTED: 2026-10-01
% ATTEMPT_STATUS: unresolved
% UnsolvedMath ID 3047; source catalogue code OPG-37222
% !TeX program = lualatex
% Research attempt by GPT-6 Astra Ultra; no claim of a new complete solution.
\documentclass[11pt,a4paper]{article}
\usepackage[margin=25mm]{geometry}
\usepackage{fontspec}
\setmainfont{lmroman10-regular.otf}[BoldFont=lmroman10-bold.otf,
 ItalicFont=lmroman10-italic.otf,BoldItalicFont=lmroman10-bolditalic.otf]
\usepackage{amsmath,amssymb,mathtools}
\usepackage{unicode-math}
\setmathfont{latinmodern-math.otf}
\AtBeginDocument{\let\setminus\smallsetminus}
\usepackage{xurl,hyperref}
\hypersetup{colorlinks=true,urlcolor=blue}
\setcounter{secnumdepth}{0}
\setlength{\parindent}{0pt}
\setlength{\parskip}{5pt}
\setlength{\emergencystretch}{3em}
\begin{document}
\section*{Dividing up the unrestricted partitions}\label{3047}
{\small UnsolvedMath ID 3047 · OPG-37222 · Combinatorics · OpenGarden}
\subsection{Definitions and mathematical statement}
Work in the ring $\mathbb Z[[x]]$ of formal power series. For each pair
of positive integers $j,m$, choose a sign $\epsilon_{j,m}\in\{-1,1\}$,
and set $\epsilon_{j,0}=1$. Form the product
\[
 F(x)=\prod_{j=1}^{\infty}\left(\sum_{m=0}^{\infty}
                      \epsilon_{j,m}x^{jm}\right)
      =\sum_{N=0}^{\infty}c_Nx^N.
\]
Each coefficient is well defined: only factors $j\le N$ and terms with
$jm\le N$ can contribute to $c_N$. No analytic convergence condition is
involved, and signs for different pairs $(j,m)$ may be chosen independently.

The question is whether one infinite choice of these signs can make
$c_N\in\{-1,0,1\}$ for every $N\ge0$. More explicitly, if a partition
of $N$ has multiplicities $(m_j)_{j\ge1}$, give it the sign
$\prod_j\epsilon_{j,m_j}$. Then
\[
 c_N=\sum_{\substack{m_j\ge0\text{ finitely supported}\\
                       \sum_jjm_j=N}}\prod_j\epsilon_{j,m_j}.
\]
Thus the same rule for a part size and its multiplicity must balance the
positive and negative partitions at every integer simultaneously. The
coefficient $c_0$ is automatically $1$. Modulo $2$, all signs are $1$, so
$c_N$ has the same parity as the ordinary partition number $p(N)$.
The signs cannot be assigned independently to individual partitions:
their multiplicative factorization through $(j,m)$ is essential. Satisfying
only a finite initial range of coefficients does not itself furnish the
required infinite assignment.
\subsection{Short English statement}
Can signs in the unrestricted-partition product be chosen so every resulting coefficient is zero, one, or minus one?
\subsection{Research attempt}
\paragraph{Scope and process.}
This is a GPT-6 Astra Ultra research attempt dated 1 October 2026, using
primary-source browsing and elementary mathematical checks. The canonical
definition above is retained. Results below are restricted results or
reductions, not a claim of a new solution to the entire catalogue question.
No formal proof assistant or external mathematical referee checked this text.


\paragraph{Literature and strategy.}
Newman's original posting [S2, R1] describes finite searches, but those
do not by themselves produce an infinite sign choice. We analyze an
extension one degree at a time, obtain an exact extension criterion,
and exhibit a valid finite prefix which has no extension. This makes
the obstacle to a naive greedy proof explicit.

\paragraph{The new variables at degree $N$.}
Assume all $\epsilon_{j,m}$ with $jm<N$ have been chosen. There are
exactly $\tau(N)$ new variables of weight $N$, namely
$\epsilon_{j,N/j}$ for divisors $j$ of $N$. A term of weight $N$
can contribute to $c_N$ only on its own, because multiplication by
any other positive-degree term exceeds degree $N$. Thus
\[
 c_N=B_N+\sum_{j\mid N}\epsilon_{j,N/j},
\]
where $B_N$ is determined entirely by earlier choices. If $t=\tau(N)$,
the new sum can be any member of $\{-t,-t+2,\ldots,t\}$. Hence an
extension with $|c_N|\le1$ exists exactly when
\[
 |B_N|\le t+1.
\]
Indeed, outside this interval every available sum has absolute value
at least two after adding $B_N$. Inside it, if zero has the required
parity it is available whenever $|B_N|\le t$; otherwise one of $1,-1$
is available. In the boundary case $|B_N|=t+1$ the parity is the
odd one and the nearest endpoint gives absolute value one.
The parity agrees with $p(N)$ automatically, since every signed
partition still contributes one modulo two.

\paragraph{A checked obstruction to arbitrary greedy extension.}
The following rows list $(\epsilon_{j,1},\ldots,
\epsilon_{j,\lfloor9/j\rfloor})$:
\[
\begin{array}{c|l}
j&\text{signs}\\ \hline
1&(1,-1,1,1,-1,1,-1,1,1)\\
2&(1,1,1,-1)\\
3&(-1,1,1)\\
4&(1,1)\\
5&(-1)\\6&(-1)\\7&(1)\\8&(-1)\\9&(1)
\end{array}
\]
Truncated multiplication gives
\[
 (c_0,c_1,\ldots,c_9)=(1,1,0,1,1,1,-1,1,0,0).
\]
All are allowed. With the weight-ten coefficients omitted, the
coefficient of $x^{10}$ is $B_{10}=6$. The four new signs are
$\epsilon_{1,10},\epsilon_{2,5},\epsilon_{5,2},\epsilon_{10,1}$,
whose sum is at least $-4$. Consequently $c_{10}\ge2$ for every
extension of this prefix. This is a counterexample to the assertion
that every valid prefix can be extended, not to Newman's existential
infinite-assignment conjecture.

\paragraph{Verification of the numerical certificate.}
The computation uses exact integers. Begin with coefficient vector
$(1,0,\ldots,0)$ through degree ten. For each row $j$, replace it
by its convolution with the factor having constant one and the
listed signed coefficients in degrees $j,2j,\ldots$. Unlisted
coefficients are zero only for this computation of $B_{10}$.
The resulting vector is
$(1,1,0,1,1,1,-1,1,0,0,6)$.
The independent verification script
\texttt{scripts/check\_attempt\_batch\_algebra\_2026\_10\_01.py}
recomputes this vector and checks all $2^4$ extensions at degree ten.
Unlisted signs in the original problem must eventually be chosen
as $\pm1$; the zero convention is just a device for isolating the
old contribution before adding the new signs explicitly.

\paragraph{What compactness would and would not give.}
At level $N$, keep all choices of signs with $jm\le N$ satisfying
the coefficient constraints through degree $N$. This is a finite
set, with finitely many extensions per node. If every level is
nonempty, there is an infinite branch: choose a child admitting
arbitrarily deep descendants at each step, which is possible by
finite branching. Every coefficient of the resulting infinite
assignment is fixed at its corresponding level, so all constraints
hold. Conversely an infinite assignment makes every level nonempty.
Thus proving nonemptiness for every degree would suffice, even
without an explicit formula for the signs. Verifying a fixed large
degree, or finding many nodes there, does not prove this premise.

\paragraph{Final status and first remaining gap.}
\textbf{UNRESOLVED.} The exact one-step criterion and a genuine
dead-end prefix are established. The missing invariant must select
prefixes that keep $|B_N|\le\tau(N)+1$ for all future degrees, or
otherwise prove arbitrarily deep feasible levels. This attempt has
no such invariant and no obstruction affecting every infinite choice.

\subsection{Sources}
\begin{itemize}
\item[S1] ULAM.AI and the UnsolvedMath Contributors, supplied problems.json, record 3047 (OPG-37222), OpenGarden collection. Catalogue identity and source transcription; CC BY 4.0.
\url{https://huggingface.co/datasets/ulamai/UnsolvedMath}.
\item[S2] Open Problem Garden, Dividing up the unrestricted partitions. Original problem and discussion; the mathematical formulation below is expanded from the supplied source record.
\url{http://www.openproblemgarden.org/op/dividing_up_the_unrestricted_partitions}.

\item[R1] David S. Newman, \emph{Dividing up the unrestricted
partitions}, original problem posting, 11 May 2010; statement and
finite-search discussion checked 1 October 2026.
\url{https://www.openproblemgarden.org/op/dividing_up_the_unrestricted_partitions}.

\end{itemize}
\end{document}
